%------------------------------------------------------------------------------%
\addtocounter{exemplo}{1}
\begin{frame}{Exemplo \theexemplo}
  Considerando a função
  \(
    f(x, y) = 4xy -x^4 -y^4
  \)

  a) Calcule o gradiente de $f$

  b) Calcule a hessiana de $f$

  c) Encontre todos os pontos críticos de $f$

  d) Classifique cada ponto crítico de $f$
\end{frame}

%------------------------------------------------------------------------------%
\begin{frame}{Exemplo \theexemplo}
  \noindent\textbf{a)}

  Precisamos das derivadas parciais de $f$
  \[
  \D{f}{x}
  = \D{}{x}\left[4xy -x^4 -y^4\right]
  = 4y - 4x^3
  \]
  \[
  \D{f}{y}
  = \D{}{y}\left[4xy -x^4 -y^4\right]
  = 4x -4y^3
  \]
  então
  \[
  \nabla f = \vetor{4y - 4x^3}{4x -4y^3}
  \]
\end{frame}

%------------------------------------------------------------------------------%
\begin{frame}{Exemplo \theexemplo}
  \noindent\textbf{b)}

  Precisamos das derivadas parciais de segunda ordem de $f$
  \[
  \DS{f}{x}
  = \D{}{x}\left[4y - 4x^3\right]
  = -12x^2
  \]
  \[
  \D{f}{y}
  = \D{}{y}\left[4x -4y^3\right]
  = -12y^2
  \]
  \[
  \DM{f}{x}{y}
  = \D{}{x}\D{f}{y}
  = \D{}{x}\left[4x -4y^3\right]
  = 4
  \]
  então
  \[
  H = \left(\begin{array}{cc}
    -12x^2 & 4 \\
    4 & -12y^2
  \end{array}\right)
  \]
\end{frame}

%------------------------------------------------------------------------------%
\begin{frame}{Exemplo \theexemplo}
  \noindent\textbf{c)}

  A função é um polinômio, então possui derivadas em todos os pontos do plano.
  Assim os pontos críticos são apenas os pontos onde as derivadas parciais são zero,
  \(\nabla f = 0\)
  \[
  4y - 4x^3 = 0
  \qquad\text{e}\qquad
  4x -4y^3 = 0
  \]
  ou, simplificando,
  \[
  y - x^3 = 0
  \qquad\text{e}\qquad
  x - y^3 = 0
  \]
  Isolando $y$ na primeira equação, \(y = x^3\), e substituindo na segunda, temos
  \begin{align*}
    x - y^3                & = 0 \\
    x - \left(x^3\right)^3 & = 0 \\
    x -x^9                 & = 0 \\
    x(1 - x^8)             & = 0
  \end{align*}
  As soluções dessa equação são $x=0$, $x=1$ ou $x-1$.
  Se $x=0$ temos $y=0$,
  se $x=1$ temos $y=1$ e
  se $x=-1$ temos $y=-1$.
  Portanto, os pontos críticos são
  \[
  (x_1, y_1) = (0, 0),
  \qquad\qquad
  (x_2, y_2) = (1, 1)
  \qquad\text{e}\qquad
  (x_3, y_3) = (-1, -1)
  \]
\end{frame}

%------------------------------------------------------------------------------%
\begin{frame}{Exemplo \theexemplo}
  \noindent\textbf{d)}

  Para classificar os pontos críticos precisamos avaliar o discriminante nos
  pontos críticos.
\end{frame}

%------------------------------------------------------------------------------%
\begin{frame}{Exemplo \theexemplo}
  Considerando o ponto \((x_1, y_1) = (0, 0)\)
  \[
  f_{xx}(0, 0) = 0
  \]
  \[
  f_{yy}(0, 0) = 0
  \]
  \[
  f_{xy}(0, 0) = 4
  \]
  \[
  D_1
  = f_{xx}(0, 0)f_{yy}(0, 0) - f^2_{xy}(0, 0)
  = 0 \times 0 - 4^2
  = -16 < 0
  \]
  Portanto, o ponto $(0, 0)$ é um ponto de sela.
\end{frame}

%------------------------------------------------------------------------------%
\begin{frame}{Exemplo \theexemplo}
  Considerando o ponto \((x_2, y_2) = (1, 1)\)
  \[
  f_{xx}(1, 1) = -12
  \]
  \[
  f_{yy}(1, 1) = -12
  \]
  \[
  f_{xy}(1, 1) = 4
  \]
  \[
  D_2
  = f_{xx}(1, 1)f_{yy}(1, 1) - f^2_{xy}(1, 1)
  = (-12)(-12) - 4^2
  = 144 -16
  = 128 > 0
  \]
  Portanto, o ponto $(1, 1)$ é um máximo ou mínimo local. Como \(f_{xx}(1, 1) = -12 < 0\)
  o ponto é um ponto de máximo local.
\end{frame}

%------------------------------------------------------------------------------%
\begin{frame}{Exemplo \theexemplo}
  Considerando o ponto \((x_3, y_3) = (-1, -1)\)
  \[
  f_{xx}(-1, -1) = -12
  \]
  \[
  f_{yy}(-1, -1) = -12
  \]
  \[
  f_{xy}(-1, -1) = 4
  \]
  \[
  D_3
  = f_{xx}(-1, -1)f_{yy}(-1, -1) - f^2_{xy}(-1, -1)
  = (-12)(-12) - 4^2
  = 144 -16
  = 128 > 0
  \]
  Portanto, o ponto $(-1, -1)$ é um máximo ou mínimo local. Como \(f_{xx}(-1, -1) = -12 < 0\)
  o ponto é um ponto de máximo local.
\end{frame}
